```latex
\documentclass[12pt]{article}

\usepackage[margin=1in]{geometry}
\usepackage{array}
\usepackage{hyperref}

\title{Cryptography and Network Security Exam}
\author{Kevine Kabarira}
\date{\today}

\begin{document}

\maketitle

\section{Introduction}

This project assesses security risks in a polytechnic institute environment and demonstrates basic cryptography, file integrity checking, and network security controls. Sample data was used throughout the project and no real student records, passwords, or encryption keys were uploaded to the repository.

\section{Risk Assessment}

\subsection{Assets, Vulnerabilities and Consequences}

\begin{center}
\begin{tabular}{|p{3cm}|p{4cm}|p{5cm}|}
\hline
\textbf{Asset} & \textbf{Vulnerability} & \textbf{Consequence} \\
\hline
Student records server & Guest network access & Unauthorised access to student records \\
\hline
Staff accounts & Weak passwords & Account compromise \\
\hline
File-transfer system & Unencrypted transfers & Data interception \\
\hline
\end{tabular}
\end{center}

\subsection{Risk Ranking}

\begin{center}
\begin{tabular}{|p{3cm}|p{2cm}|p{2cm}|p{7cm}|}
\hline
\textbf{Risk} & \textbf{Likelihood} & \textbf{Impact} & \textbf{Reason} \\
\hline
Guest access & High & High & Guests can reach sensitive records. \\
\hline
Weak passwords & High & High & Weak passwords can be guessed or compromised. \\
\hline
Unencrypted transfers & Medium & High & Data can be intercepted during transfer. \\
\hline
\end{tabular}
\end{center}

\subsection{Recommended Controls}

\begin{itemize}
\item Guest access: Use firewall rules and network segmentation.
\item Weak passwords: Enforce a strong password policy.
\item Unencrypted transfers: Use SFTP.
\end{itemize}

\section{Python Security Toolkit}

A Python security toolkit was developed to demonstrate encryption, decryption, hashing, integrity verification, and error handling.

\subsection{Encryption}

The toolkit uses Fernet symmetric encryption from the Python \texttt{cryptography} library. A sample student record was encrypted and saved as an encrypted output file.

The encryption key was stored outside the GitHub repository to prevent exposure.

\subsection{Decryption and Verification}

The encrypted file was successfully decrypted. The SHA-256 hash of the original file was compared with the SHA-256 hash of the decrypted file.

The results were:

\begin{verbatim}
Original SHA-256:
9072923fe0a1c93d2fe2f0f564281c78c0d6e7cb9ddb1a315578fcbaf428b877

Decrypted SHA-256:
9072923fe0a1c93d2fe2f0f564281c78c0d6e7cb9ddb1a315578fcbaf428b877

Contents match: True
Integrity check: PASS
\end{verbatim}

The matching hashes confirm that the decrypted file has the same contents as the original sample file.

\subsection{Integrity Detection}

A copy of the sample file was modified after the original hash was calculated. The hashes were:

\begin{verbatim}
Original hash:
9072923fe0a1c93d2fe2f0f564281c78c0d6e7cb9ddb1a315578fcbaf428b877

Modified file hash:
50d29235c461aec22174106942a54ea16096fbb5a55aacdb52ba74150cdb9f

Modified file detected: True
\end{verbatim}

This demonstrates that SHA-256 can detect changes to the file.

\subsection{Error Handling}

The tool
```
